%----------------------------------------------------------------------------------------
%   USEFUL COMMANDS
%----------------------------------------------------------------------------------------

\newcommand{\dipartimento}{Dipartimento di Matematica ``Tullio Levi-Civita''}

%----------------------------------------------------------------------------------------
% 	USER DATA
%----------------------------------------------------------------------------------------

% Data di approvazione del piano da parte del tutor interno; nel formato GG Mese AAAA
% compilare inserendo al posto di GG 2 cifre per il giorno, e al posto di 
% AAAA 4 cifre per l'anno
\newcommand{\dataApprovazione}{Data}

% Dati dello Studente
\newcommand{\nomeStudente}{Gabriele}
\newcommand{\cognomeStudente}{Scaggiante}
\newcommand{\matricolaStudente}{2101076}
\newcommand{\emailStudente}{gabriele.scaggiante.1@studenti.unipd.it}
\newcommand{\telStudente}{+ 39 375 583 0717}

% Dati del Tutor Aziendale
\newcommand{\nomeTutorAziendale}{Fabio}
\newcommand{\cognomeTutorAziendale}{Pallaro}
\newcommand{\emailTutorAziendale}{f.pallaro@synclab.it}
\newcommand{\telTutorAziendale}{+ 39 333 136 8500}
\newcommand{\ruoloTutorAziendale}{}

% Dati dell'Azienda
\newcommand{\ragioneSocAzienda}{Sync Lab S.r.l.}
\newcommand{\indirizzoAzienda}{Galleria Spagna, 28, 35127 Padova}
\newcommand{\sitoAzienda}{https://synclab.it/}
\newcommand{\emailAzienda}{info@synclab.it}
\newcommand{\partitaIVAAzienda}{P.IVA 07952560634}

% Dati del Tutor Interno (Docente)
\newcommand{\titoloTutorInterno}{Prof.}
\newcommand{\nomeTutorInterno}{NomeDocente}
\newcommand{\cognomeTutorInterno}{CognomeDocente}

\newcommand{\prospettoSettimanale}{
     % Personalizzare indicando in lista, i vari task settimana per settimana
     % sostituire a XX il totale ore della settimana
    \begin{itemize}
        \item \textbf{Prima Settimana (24 ore)}
        \begin{itemize}
            \item Introduzione ad Angular
        \end{itemize}
        \item \textbf{Seconda Settimana (24 ore)} 
        \begin{itemize}
            \item Studio approfondito di Angular
            \item Introduzione a Java e Spring
        \end{itemize}
        \item \textbf{Terza Settimana (24 ore)} 
        \begin{itemize}
            \item Studio approfondito di Spring
        \end{itemize}
        \item \textbf{Quarta Settimana (40 ore)} 
        \begin{itemize}
            \item Studio completo dei requisiti
            \item Studio completo dei casi d'uso
            \item Progettazione dell'architettura dell'applicazione
            \item Progettazione dello schema del database
        \end{itemize}
        \item \textbf{Quinta Settimana (40 ore)} 
        \begin{itemize}
            \item Sviluppo frontend della schermata di Homepage
            \item Sviluppo frontend della schermata di registrazione
            \item Sviluppo frontend della bacheca dei post
            \item Sviluppo della funzionalità di autenticazione e di gestione dell'account utente
        \end{itemize}
        \item \textbf{Sesta Settimana (40 ore)} 
        \begin{itemize}
            \item Sviluppo frontend del singolo post
            \item Sviluppo backend delle operazioni CRUD per i post
        \end{itemize}
        \item \textbf{Settima Settimana (40 ore)} 
        \begin{itemize}
            \item Sviluppo delle funzionalità di commento e revisione ai post (backend e frontend)
        \end{itemize}
        \item \textbf{Ottava Settimana (40 ore)} 
        \begin{itemize}
            \item Integrazione completa tra frontend e backend
            \item Gestione degli errori e validazione dei dati (frontend e backend)
            \item Implementazione dei test (unit test e test di integrazione)
            \item Refactoring del codice e ottimizzazione delle prestazioni
        \end{itemize}
        \item \textbf{Nona Settimana (40 ore)} 
        \begin{itemize}
            \item Verifica conformità ai requisiti individuati
            \item Rifinitura dell'applicazione e miglioramenti dell'interfaccia utente
            \item Verifica e miglioramento degli aspetti di sicurezza dell'applicazione
        \end{itemize}
    \end{itemize}
}

% Indicare il totale complessivo (deve essere compreso tra le 300 e le 320 ore)
\newcommand{\totaleOre}{312}

\newcommand{\obiettiviObbligatori}{
     \item \underline{\textit{O01}}: Studio e comprensione delle tecnologie di frontend utilizzate;
     \item \underline{\textit{O02}}: Studio e comprensione delle tecnologie di backend utilizzate;
     \item \underline{\textit{O03}}: Studio e comprensione del problema da risolvere, dell'architettura del software e del suo funzionamento;
     \item \underline{\textit{O04}}: Progettazione dell'architettura del sistema e definizione dei componenti software;
	 \item \underline{\textit{O05}}: Sviluppo della funzionalità di autenticazione all'applicazione;
	 \item \underline{\textit{O06}}: Sviluppo della schermata introduttiva della piattaforma;
	 \item \underline{\textit{O07}}: Sviluppo di una schermata interattiva contenente i post pubblicati dagli utenti registrati alla piattaforma;
	 \item \underline{\textit{O08}}: Sviluppo della schermata di dettaglio relativa a un singolo post. La schermata dovrà mostrare i contenuti del post originale, l'autore del post e le interazioni da parte degli altri utenti iscritti alla piattaforma;
	 \item \underline{\textit{O09}} Sviluppo della funzionalità CRUD per i post;
	 \item \underline{\textit{O10}} Sviluppo della funzionalità CRUD per i commenti ai post;
	 \item \underline{\textit{O11}}: Progettazione e sviluppo delle API per la comunicazione tra frontend e backend;
	 \item \underline{\textit{O12}}: Integrazione tra frontend e backend e gestione della comunicazione client-server;
	 \item \underline{\textit{O13}} Verifica della conformità del prodotto rispetto agli obiettivi.
}

\newcommand{\obiettiviDesiderabili}{
    \item \underline{\textit{D01}}: Containerizzazione dell'applicazione tramite Docker;
    \item \underline{\textit{D02}}: Implementazione di test di unità e di integrazione;
    \item \underline{\textit{D03}}: Sviluppo di una schermata di gestione dell'account utente;
    \item \underline{\textit{D04}}: Visualizzazione di informazioni di dettaglio sugli utenti all'interno dei post e delle interazioni;
}

\newcommand{\obiettiviFacoltativi}{
	 \item \underline{\textit{F01}}: Realizzazione di una funzionalità di chat in tempo reale all'interno dei post;
	 \item \underline{\textit{F02}}: Esecuzione degli snippet di codice presenti all'interno dei post;
	 \item \underline{\textit{F03}}: Implementazione di un sistema di valutazione dei post e dei commenti;
	 \item \underline{\textit{F04}}: Introduzione di un sistema di notifiche per aggiornamenti e interazioni sui post;
	 \item \underline{\textit{F05}}: Implementazione di funzionalità di ricerca e filtraggio dei post;
	 \item \underline{\textit{F06}}: Implementazione di un sistema di tag per l'organizzazione dei contenuti;
}